% Part: incompleteness
% Chapter: theories-computability
% Section: consis-with-q

\documentclass[../../../include/open-logic-section]{subfiles}

\begin{document}

\olfileid[hi]{inc}{tcp}{con}

\olsection{$\Th{Q}$ के साथ संगत सिद्धांत अनिर्णेय हैं}

अगला प्रमेय कहता है कि केवल $\Th{Q}$ ही अनिर्णेय नहीं है, बल्कि वास्तव में
$\Th{Q}$ से असहमत न होने वाला प्रत्येक सिद्धांत अनिर्णेय है।

\begin{thm}
$\Th{T}$ अंकगणित की भाषा का ऐसा कोई सिद्धांत हो जो $\Th{Q}$ के साथ संगत
है (अर्थात् $\Th{T} \cup \Th{Q}$ संगत है)। तब $\Th{T}$ अनिर्णेय है।
\end{thm}

\begin{proof}
याद करें कि $\Th{Q}$ के अभिगृहीतों का परिमित समुच्चय $!Q_1$, \dots,
$!Q_8$ है। हम इनके स्थान पर एक ही अभिगृहीत
$!E = !Q_1 \land \dots \land !Q_8$ भी रख सकते हैं।

मान लें $\Th{T}$ कोई निर्णेय सिद्धांत है जो~$\Th{Q}$ के साथ संगत है। मानें
\[
C = \Setabs{!A}{\Th{T} \Proves !E \lif !A}.
\]
हम दिखाते हैं कि $C$,~$\Th{Q}$ और $\Th{\bar Q}$ का संगणनीय वियोजन होगा,
जो विरोधाभास है। पहले, यदि $!A$,~$\Th{Q}$ में है, तो $!A$,~$\Th{Q}$ के
अभिगृहीतों से सिद्ध है; निगमन प्रमेय से प्रथम-कोटि तर्कशास्त्र में
$!E \lif !A$ की कोई !!a{derivation} है। अतः $!A$,~$C$ में है।

दूसरी ओर यदि $!A$,~$\Th{\bar Q}$ में है, तो प्रथम-कोटि तर्कशास्त्र में
$!E \lif \lnot !A$ का प्रमाण है। यदि $\Th{T}$,~$!E \lif !A$ को भी सिद्ध
करता है, तो $\Th{T}$,~$\lnot !E$ को सिद्ध करता है; उस स्थिति में
$\Th{T} \cup \Th{Q}$ असंगत है। किंतु हम मान रहे हैं कि
$\Th{T} \cup \Th{Q}$ संगत है, इसलिए $\Th{T}$,~$!E \lif !A$ को सिद्ध
नहीं करता और इस प्रकार $!A$,~$C$ में नहीं है।

हमने दिखाया है कि यदि $!A$,~$\Th{Q}$ में है, तो वह $C$ में है; और यदि
$!A$,~$\Th{\bar Q}$ में है, तो वह $\Complement{C}$ में है। अतः $C$ एक
संगणनीय वियोजन है, और यही हमारा वांछित विरोधाभास है।
\end{proof}

यह प्रमेय बहुत शक्तिशाली है। उदाहरणतः, इससे निकलता है:

\begin{cor}
  अंकगणित की भाषा के लिए प्रथम-कोटि तर्कशास्त्र (अर्थात् समुच्चय
  $\Setabs{!A}{\text{$!A$ प्रथम-कोटि तर्कशास्त्र में सिद्ध है}}$)
  अनिर्णेय है।
\end{cor}

\begin{proof}
प्रथम-कोटि तर्कशास्त्र $\emptyset$ के निष्पादों का समुच्चय है, और
वह~$\Th{Q}$ के साथ संगत है।
\end{proof}

\end{document}
